\section{附录：符号与术语表}\label{sec:appendix}

\begin{table}[htbp]
    \centering
    \caption{本文使用的主要符号}
    \label{tab:notation}
    \begin{tabular}{ll}
        \toprule
        符号 & 含义 \\
        \midrule
        $\zeta(s)$ & 黎曼 $\zeta$ 函数 \\
        $\pi(x)$ & 素数计数函数，$\#\{p \le x: p \text{ 为素数}\}$ \\
        $\operatorname{Li}(x)$ & 对数积分 $\int_2^x dt/\log t$ \\
        $\Lambda(n)$ & von Mangoldt 函数 \\
        $\psi(x)$ & 切比雪夫函数 $\sum_{n \le x} \Lambda(n)$ \\
        $\sigma(n)$ & 因子和函数 $\sum_{d \mid n} d$ \\
        $H_n$ & 调和数 $\sum_{k=1}^{n} 1/k$ \\
        $\gamma$ & 欧拉—马歇罗尼常数 \\
        $\rho$ & $\zeta(s)$ 的非平凡零点 \\
        $N(T)$ & $0 < \Im(\rho) \le T$ 的零点个数 \\
        $N_0(T)$ & 临界线上 $0 < \Im(\rho) \le T$ 的零点个数 \\
        $\kappa(T)$ & 零点比例 $N_0(T)/N(T)$ \\
        $\Theta$ & 零点实部的上确界 $\sup\{\Re(\rho)\}$ \\
        GUE & 高斯酉系综 \\
        GRH & 广义黎曼假设 \\
        \bottomrule
    \end{tabular}
\end{table}

\section*{术语对照}
\begin{itemize}
    \item 临界线（critical line）：复平面上的直线 $\Re(s) = \tfrac{1}{2}$；
    \item 临界带（critical strip）：竖直区域 $0 < \Re(s) < 1$；
    \item 显式公式（explicit formula）：联系素数计数与零点求和的恒等式；
    \item 硬性定理（Hardy-type theorem）：关于临界线上零点存在性或比例的无条件结果；
    \item 次凸性（subconvexity）：突破凸性界的 $L$ 函数上界估计。
\end{itemize}
